\documentclass[preview]{standalone}

\usepackage{amsmath}
\usepackage{amssymb}
\usepackage{stellar}
\usepackage{enumitem}
\usepackage{definitions}

\begin{document}

\id{linearalgebra-exercises-batch-15}
\genpage

\section{Exercises}

\begin{snippetexercise}{linear-algebra-batch-15-ex-1}{}
    Classify metrically the conics defined by the following equations:
    \begin{enumerate}[label=\Alph*]
        \item \(x^2 - y^2 + 2xy + 2x - \sqrt{3} = 0\);
        \item \(5x^2 + 7y^2 - 4xy + 4x + 3y - \sqrt{11} = 0\);
        \item \(x^2 + y^2 + 2xy + x + y = 0\);
        \item \(2x^2 + 5y^2 + 6xy + 4x + 10y + 10 = 0\);
        \item \(37x^2 + 13y^2 - 18xy + 92x - 44y + 28 = 0\).
    \end{enumerate}
\end{snippetexercise}

\begin{snippetsolution}{linear-algebra-batch-15-ex-1-sol}{}
    \todo
\end{snippetsolution}

\begin{snippetexercise}{linear-algebra-batch-15-ex-2}{}
    Determine for each \(k \in \realnumbers\) the affine type
    of the conics defined by the following equations:
    \begin{enumerate}[label=\Alph*]
        \item \(5x^2 + 5y^2 - 6xy - 10x - k = 0\);
        \item \((k+1)x^2 + (k-1)y^2 + 2kx -y - 1 = 0\).
    \end{enumerate}
\end{snippetexercise}

\begin{snippetsolution}{linear-algebra-batch-15-ex-2-sol}{}
    \todo
\end{snippetsolution}

\begin{snippetexercise}{linear-algebra-batch-15-ex-3}{}
    After verifying that the conic
    \[
        \mathcal C \triangleq \{
            (x,y) \in \mathcal E^2 \suchthat 10x^2 - 15y^2 + 24xy - 4x + 54y - 23 = 0
        \}
    \]
    is a hyperbola, determine an affine isometry of \(\mathcal E^2\) which transforms
    \(\mathcal C\) into an hyperbola of equation
    \[
        \frac{X^2}{a^2} - \frac{Y^2}{b^2} = 1
    \]
\end{snippetexercise}

\begin{snippetsolution}{linear-algebra-batch-15-ex-3-sol}{}
    \todo
\end{snippetsolution}

\begin{snippetexercise}{linear-algebra-batch-15-ex-4}{}
    After verifying that the conic
    \[
        \mathcal C \triangleq \{
            (x,y) \in \mathcal E^2 \suchthat 10x^2 - 15y^2 + 24xy - 4x + 54y - 23 = 0
        \}
    \]
    is a hyperbola, determine an affine isometry of \(\mathcal E^2\) which transforms
    \(\mathcal C\) into an hyperbola of equation
    \[
        \frac{X^2}{a^2} - \frac{Y^2}{b^2} = 1
    \]
\end{snippetexercise}

\begin{snippetsolution}{linear-algebra-batch-15-ex-4-sol}{}
    \todo
\end{snippetsolution}

\begin{snippetexercise}{linear-algebra-batch-15-ex-5}{}
    Let \(\varphi \colon \mathcal V \fromto \mathcal W\) be an affine map.
    \begin{enumerate}[label=\Alph*]
        \item Show that \(\varphi\) is completely
        determined by the image of an affine frame.
        \item Let \(\mathcal V = \mathcal W = \mathcal A^2(\realnumbers)\) and
        \(\varphi\) the only affine map such that \((0,0) \fromto (-1, 0), (1,1) \fromto (-2, -1)\)
        and \((2,1) \fromto (-3, 1)\).
        Find the \matrix associated to \(\varphi\)
        in the affine frame \(\{0, e_1, e_2\}\).
    \end{enumerate}
\end{snippetexercise}

\begin{snippetsolution}{linear-algebra-batch-15-ex-5-sol}{}
    \todo
\end{snippetsolution}

\begin{snippetexercise}{linear-algebra-batch-15-ex-6}{}
    In \(\mathcal E^3\) determine the equation of the line passing through
    \(p = (1,1,1)\) and
    \begin{enumerate}[label=\Alph*]
        \item for \(q = (2,3,-1)\);
        \item is parallel to the line
        \[
            \begin{cases}
                x = -1 + t \\
                y =1-t \\
                z = 2-3t
            \end{cases}
        \]
        \item is perpendicular to the plane containing the lines
        \[
            \begin{cases}
                x = 2 + t \\
                y = 1 + 2t \\
                z = 2 + 9t
            \end{cases}, \quad
            \begin{cases}
                x = 1 + z \\
                y = 1 - u \\
                z = 1 - 3u
            \end{cases}
        \]
    \end{enumerate}
\end{snippetexercise}

\begin{snippetsolution}{linear-algebra-batch-15-ex-6-sol}{}
    \todo
\end{snippetsolution}

\end{document}